\documentclass[11pt,a4paper]{article}

\usepackage[utf8]{inputenc}
\usepackage[T1]{fontenc}
\usepackage{lmodern}
\usepackage[margin=0.76in]{geometry}
\usepackage[table]{xcolor}
\usepackage{array}
\usepackage{tabularx}
\usepackage{booktabs}
\usepackage{enumitem}
\usepackage{titlesec}
\usepackage{fancyhdr}
\usepackage{lastpage}
\usepackage{hyperref}
\usepackage[normalem]{ulem}

\definecolor{TextInk}{HTML}{233142}
\definecolor{HeadingBlue}{HTML}{1D3557}
\definecolor{AccentTeal}{HTML}{2A7F92}
\definecolor{AccentGold}{HTML}{C58B4E}
\definecolor{RuleGray}{HTML}{D6E0E5}
\definecolor{TableStripe}{HTML}{F3F8FA}
\definecolor{SoftTint}{HTML}{F7F4EE}
\definecolor{SoftBlue}{HTML}{EEF6F8}
\definecolor{SoftGreen}{HTML}{EEF7F0}
\definecolor{SoftRed}{HTML}{FBF0EF}

\hypersetup{
  colorlinks=true,
  linkcolor=HeadingBlue,
  urlcolor=AccentTeal,
  pdftitle={IB Geography SL May 5 Food Health And Climate Visual Remediation},
  pdfauthor={IB Geography SL Preparation}
}

\color{TextInk}
\setlength{\parskip}{0.45em}
\setlength{\parindent}{0pt}
\setlength{\tabcolsep}{5pt}
\renewcommand{\arraystretch}{1.18}
\setlist[itemize]{leftmargin=1.32em,itemsep=3pt,topsep=4pt}
\setlist[enumerate]{leftmargin=1.45em,itemsep=3pt,topsep=4pt}

\newcolumntype{Y}{>{\raggedright\arraybackslash}X}
\newcolumntype{L}[1]{>{\raggedright\arraybackslash}p{#1}}
\newcolumntype{C}[1]{>{\centering\arraybackslash}p{#1}}

\titleformat{\section}
  {\normalfont\huge\sffamily\bfseries\color{HeadingBlue}}
  {}{0pt}{}
  [\vspace{0.12em}{\color{AccentGold}\titlerule[1pt]}]
\titleformat{\subsection}
  {\normalfont\Large\sffamily\bfseries\color{AccentTeal}}
  {}{0pt}{}
\titleformat{\subsubsection}
  {\normalfont\large\sffamily\bfseries\color{HeadingBlue}}
  {}{0pt}{}
\titlespacing*{\section}{0pt}{1.05em}{0.45em}
\titlespacing*{\subsection}{0pt}{0.78em}{0.22em}
\titlespacing*{\subsubsection}{0pt}{0.55em}{0.16em}

\pagestyle{fancy}
\setlength{\headheight}{14pt}
\fancyhf{}
\fancyhead[L]{\sffamily\color{AccentTeal}Tue 2026-05-05}
\fancyhead[R]{\sffamily\color{HeadingBlue}IB Geography SL}
\fancyfoot[C]{\sffamily\color{HeadingBlue}\thepage\ of \pageref{LastPage}}
\renewcommand{\headrulewidth}{0.4pt}

\newcommand{\term}[1]{\textcolor{HeadingBlue}{\textbf{#1}}}
\newcommand{\exam}[1]{\textcolor{AccentTeal}{\textbf{#1}}}
\newcommand{\todobox}{\fbox{\rule{0pt}{0.8em}\rule{0.8em}{0pt}}}
\newcommand{\boxnote}[3]{%
\noindent\colorbox{#1}{%
  \begin{minipage}{0.965\textwidth}
  \sffamily
  \textbf{\textcolor{HeadingBlue}{#2}} #3
  \end{minipage}
}}

\begin{document}

\begin{center}
  {\sffamily\bfseries\color{HeadingBlue}\LARGE IB Geography SL Study Pack}\\[0.25em]
  {\sffamily\bfseries\color{AccentTeal}\Large Tuesday 2026-05-05}\\[0.35em]
  {\sffamily\color{TextInk}Food And Health Remediation + Climate Visual-Stimulus Practice}
\end{center}

\boxnote{SoftBlue}{How this pack follows the plan.}{The May 5 schedule says:
\term{Audio 1}: weakest \term{Food and health} content only.
\term{Audio 2}: \term{Climate} infographic and policy evaluation practice.
\term{25 min}: do one fast visual-stimulus response set. From today onward,
the preparation plan says to stop expanding notes and focus on timed repair.}

\section{Today's Target}

\rowcolors{2}{TableStripe}{white}
\begin{tabularx}{\textwidth}{L{0.2\textwidth}Y L{0.25\textwidth}}
\toprule
\textbf{Part} & \textbf{Focus} & \textbf{Output} \\
\midrule
Audio 1 & Fix weakest Food and health areas: source evidence, double burden,
stakeholders, strategy evaluation & One Food and health repair checklist \\
Audio 2 & Climate infographic reading and policy evaluation & One visual-source
routine and one policy-fit sentence \\
25 min block & One fast visual-stimulus response set & One 10-mark response
set plus marking and error-log sentence \\
\bottomrule
\end{tabularx}
\rowcolors{2}{}{}

\boxnote{SoftRed}{No new-content drift.}{Use only the existing Food and health
and Climate habits from the April 26, April 30, May 2, and May 3 materials.
The aim is faster source use and clearer evaluation, not more case studies.}

\section{Audio 1 Notes: Food And Health Repair}

\subsection{Weakness 1: Source-Led Food And Health Contrasts}

The May 2 Food and health marking guide rewards exact source contrasts. Use:

\begin{center}
{\sffamily\bfseries\color{HeadingBlue}
Region + indicator + number + comparison + brief meaning}
\end{center}

\rowcolors{2}{TableStripe}{white}
\begin{tabularx}{\textwidth}{L{0.24\textwidth}Y Y}
\toprule
\textbf{Question type} & \textbf{Risky answer} & \textbf{Better answer} \\
\midrule
\exam{State} contrast & Region X is less healthy. & Region X has
\term{38\%} food insecurity, compared with \term{7\%} in Region Z. \\
\exam{Describe} change & Health problems change. & Child undernutrition falls
from \term{28\%} in Region X to \term{3\%} in Region Z, while adult obesity
rises from \term{6\%} to \term{32\%}. \\
\exam{Explain} stakeholder effect & TNCs affect food. & TNCs can improve
distribution and jobs, but may worsen health if marketing and pricing
encourage ultra-processed diets. \\
\bottomrule
\end{tabularx}
\rowcolors{2}{}{}

\subsection{Food Security And Health Toolkit}

\rowcolors{2}{TableStripe}{white}
\begin{tabularx}{\textwidth}{L{0.22\textwidth}Y Y}
\toprule
\textbf{Idea} & \textbf{Use in answers} & \textbf{Evaluation warning} \\
\midrule
Availability & Is there enough food in the place? & Enough national supply
does not guarantee household access. \\
Access & Can people afford and physically reach food? & Access depends on
income, price, transport, markets, and distribution. \\
Utilization & Is food safe, nutritious, and usable by the body? & Clean water,
sanitation, disease, and diet quality matter. \\
Stability & Does access continue through shocks? & Drought, conflict, price
shocks, illness, and supply disruption can reverse gains. \\
Double burden & Undernutrition and obesity can coexist. & Do not assume one
region has only one health problem. \\
\bottomrule
\end{tabularx}
\rowcolors{2}{}{}

\subsection{Stakeholder Chain}

\boxnote{SoftGreen}{Repair formula.}{Stakeholder $\rightarrow$ action
$\rightarrow$ food-security dimension $\rightarrow$ health outcome
$\rightarrow$ limitation.}

\begin{itemize}
  \item \term{Government}: school meals, safety nets, water and sanitation,
  public health, farmer support; limited by cost, targeting, enforcement, and
  political priorities.
  \item \term{TNC}: investment, logistics, retail, jobs, processing, marketing;
  limited by market power, processed-food promotion, nutrition quality, and
  regulation.
  \item \term{NGO / international organization}: emergency food, health
  education, vaccination, clean water, clinics; limited by donor dependence and
  short project timescales.
  \item \term{Households and consumers}: diet, waste, health behaviour; choices
  are constrained by income, time, culture, transport, and advertising.
\end{itemize}

\section{Audio 2 Notes: Climate Visuals}

\subsection{Visual-Stimulus Routine}

\rowcolors{2}{TableStripe}{white}
\begin{tabularx}{\textwidth}{L{0.18\textwidth}Y}
\toprule
\textbf{Step} & \textbf{Action} \\
\midrule
Read & Title, categories, units, index labels, legend, and command terms. \\
Describe & Highest, lowest, clear contrast, and one exception or nuance. \\
Explain & Link hazard, exposure, sensitivity, and capacity to respond. \\
Evaluate & Match the policy to the risk, then give one limitation. \\
\bottomrule
\end{tabularx}
\rowcolors{2}{}{}

\subsection{Policy-Fit Criteria}

\rowcolors{2}{TableStripe}{white}
\begin{tabularx}{\textwidth}{L{0.2\textwidth}Y Y}
\toprule
\textbf{Criterion} & \textbf{Question to ask} & \textbf{Useful language} \\
\midrule
Effectiveness & Does it reduce the risk shown by the source? & This reduces
exposure / sensitivity by... \\
Equity & Does it reach the most vulnerable group? & It is stronger if it
protects informal housing / low-income households / outdoor workers. \\
Feasibility & Can it be funded, maintained, and governed? & It may be limited
by cost, maintenance, training, or political capacity. \\
Sustainability & Does it reduce long-term risk? & It is only partly successful
if it shifts risk or ignores emissions / land-use planning. \\
\bottomrule
\end{tabularx}
\rowcolors{2}{}{}

\boxnote{SoftTint}{Climate answer chain.}{Hazard $\rightarrow$ exposed group
$\rightarrow$ vulnerability factor $\rightarrow$ policy response
$\rightarrow$ benefit $\rightarrow$ limitation.}

\section{25-Minute Visual-Stimulus Block}

\rowcolors{2}{TableStripe}{white}
\begin{tabularx}{\textwidth}{L{0.16\textwidth}Y L{0.24\textwidth}}
\toprule
\textbf{Time} & \textbf{Task} & \textbf{Output} \\
\midrule
0--3 min & Read the source title, indicators, units, and command terms. Mark
highest, lowest, and one contrast. & Annotated source \\
3--15 min & Complete the workbook's 10-mark visual-stimulus response set. &
Timed response \\
15--21 min & Mark with the May 5 guide. Identify one source-use issue and one
policy-evaluation issue. & Score and reason \\
21--25 min & Write the final error-log sentence. & One habit for Paper 2 \\
\bottomrule
\end{tabularx}
\rowcolors{2}{}{}

\section{Carry-Forward}

\begin{itemize}
  \item For \term{Food and health}, carry forward source precision, the four
  food-security dimensions, the double burden, and stakeholder chains.
  \item For \term{Climate}, carry forward visual-source discipline and policy
  fit: risk shown by source, response, benefit, limitation.
  \item For the final 72 hours before Paper 1, do not expand notes. Fix only
  the answer habits that cost marks.
\end{itemize}

\end{document}
